\section{Introduction}
\label{sec:introduction}

Capacitance loss and equivalent series resistance (ESR) growth alter electrical stress, thermal loss, and filtering performance in power converters. Online monitoring methods therefore span ripple analysis, injected excitation, circuit reconfiguration, observers, recursive identification, and data-driven processing \cite{dang2020review,nathan2023review}. Their applicability depends on the capacitor's circuit role and on which electrical variables are already measured.

The energy-transfer capacitor in a \Cuk{} converter is distinct from a dc-link or output filter: it lies in the main power-transfer path and alternately carries $+i_{L1}$ and $-i_{L2}$. Normal switching thus creates a bidirectional current commutation that can expose capacitor health without an added diagnostic perturbation. Treating this component as a generic two-terminal capacitor sampled asynchronously discards that topology-specific information.

Prior joint C--ESR methods use inherent ripple or within-cycle relations \cite{hasegawa2018dclink,yao2015buck,tang2019sensorless}, switched networks and observers \cite{barbulescu2022observer,barbulescu2023implementation,caiman2024quasionline}, or inherent and injected spectral features \cite{ribeiro2025inherent,ribeiro2024injection,ribeiro2023wavelet}. Related recent work includes sensorless Boost ESR estimation, optimized Buck-capacitor injection, and large-signal C--ESR estimation \cite{leong2025boostesr,zhang2025injection,zhang2025largesignal}. Wavelet/Kalman and recursive RLS/Kalman formulations provide further signal-processing and estimation routes \cite{amaral2024waveletkf,zhang2025haar,reichbach2016rls,chen2018vffrls}. \Cuk{}-specific modeling, fault diagnosis, and prognostics have also been studied \cite{wang2021fractionalcuk,prathiba2023cuk,ajra2024sensorfault,samie2015prognostic}. These studies establish the broader tools; the gap addressed here is how the \Cuk{} transfer mechanism itself can be separated into implementable edge and charge information.

This paper reconstructs the switching-edge discontinuity for an ESR-dominant observation and uses an ESR-compensated in-state charge balance for a capacitance-dominant observation. Table~\ref{tab:related_work} positions that scope against representative methods. The comparison is categorical because the converter topologies, sensing sets, and reporting protocols differ.

\begin{table*}[t]
\caption{Representative Related Methods and the Scope of This Work}
\label{tab:related_work}
\centering
\footnotesize
\setlength{\tabcolsep}{3pt}
\begin{tabular}{p{2.15cm}p{1.45cm}p{1.0cm}p{1.55cm}p{2.45cm}p{1.95cm}p{4.55cm}}
\toprule
Method & Capacitor role & Joint C/ESR & Added excitation & Measurements & Processing & Distinction from this work \\
\midrule
Hasegawa \emph{et al.} \cite{hasegawa2018dclink} & inverter dc link & yes & no & dc-link voltage and ripple current & ripple-frequency separation & different topology and ripple bands \\ \addlinespace[1pt]
Yao \emph{et al.} \cite{yao2015buck} & Buck output filter & yes & no & PWM timing and output voltage & within-cycle relation & topology-specific relation without \Cuk{} current reversal \\ \addlinespace[1pt]
Barbulescu \emph{et al.} \cite{barbulescu2022observer} & device under test & yes & switched network & terminal voltage and switched resistive network & parameter observer & requires an additional discharge network \\ \addlinespace[1pt]
Ribeiro \emph{et al.} \cite{ribeiro2025inherent} & NPC dc link & yes & no & converter voltage/current signals & spectral features plus RLS & different inherent-signal mapping \\ \addlinespace[1pt]
Amaral \emph{et al.} \cite{amaral2024waveletkf} & boost output/dc link & yes & no & converter signals & wavelet plus Kalman filtering & reconstructed frequency features rather than edge/charge rows \\ \addlinespace[1pt]
This work & \Cuk{} energy transfer & yes & no & $v_T$, $i_1$, $i_2$, PWM state/time & edge fit and charge balance & uses $+i_{L1}\leftrightarrow-i_{L2}$ commutation \\
\bottomrule
\end{tabular}
\vspace{1mm}

\begin{minipage}{0.98\textwidth}
\scriptsize\emph{Sensing note:} No diagnostic injection does not imply zero sensing burden. The device-realistic realization uses separate absolute- and edge-voltage ranges in the differential $v_T$ analog front end, but no capacitor-branch current sensor.
\end{minipage}
\end{table*}

The contributions are fourfold.
\begin{enumerate}
\item The normal \Cuk{} switching sequence is interpreted as topology-induced health excitation through the $+i_{L1}\leftrightarrow-i_{L2}$ capacitor-current commutation.
\item Timestamp-reconstructed edge and safe-window charge observations provide ESR-dominant and capacitance-dominant information while preserving the distinction between terminal and internal capacitor voltage.
\item Positive finite-window information bounds establish local identifiability under valid continuous-conduction-mode conditions. \TSRLS{} is the primary low-cost realization, while \TSKF{} is an uncertainty-aware extension; the standard RLS and Kalman recursions are not claimed as novelties.
\item A common evaluation protocol combines independent plant traces, 48 blind cases, noise/timing sweeps, an observation--estimator factorial, degradation trajectories, operating transients, and a TMS320F28379D device-realistic simulation.
\end{enumerate}

The remainder of the paper develops the model and observations, proves the finite-window bounds, presents the two realizations, and evaluates their accuracy, dynamics, uncertainty, and embedded feasibility.
